% ==============================================================================
% EJERCICIOS RESUELTOS - ECONOMÍA ESPAÑOLA
% ==============================================================================

\section{Cuaderno 1: Análisis de la Economía Española}

\subsection{Ejercicio 1: Estructura del PIB y Crecimiento Nominal}

\textbf{Enunciado:} Con los datos del Cuadro 1, calcular la estructura porcentual del PIB desde la demanda y por sectores en 2008, 2019 y 2023; la tasa de crecimiento interanual (2020-2023) y media (2008-2023); el PIB en números índices (base 2020); e identificar las fases económicas.

\subsubsection{1. Estructura porcentual del PIB desde la demanda}
El cálculo se realiza dividiendo cada componente entre el PIB a precios de mercado de su respectivo año.

\begin{itemize}
    \item \textbf{2008:} Consumo final (76,25\%), FBC (28,48\%), Exportaciones (25,55\%), Importaciones (-30,28\%).
    \item \textbf{2019:} Consumo final (76,11\%), FBC (20,91\%), Exportaciones (34,68\%), Importaciones (-31,69\%).
    \item \textbf{2023:} Consumo final (75,07\%), FBC (21,00\%), Exportaciones (38,06\%), Importaciones (-34,13\%).
\end{itemize}

\subsubsection{2. Estructura porcentual del PIB por sectores económicos (Oferta)}
\begin{itemize}
    \item \textbf{2008:} Agricultura (2,35\%), Industria (15,53\%), Construcción (10,66\%), Servicios (63,63\%), Impuestos netos (7,83\%).
    \item \textbf{2019:} Agricultura (2,52\%), Industria (14,09\%), Construcción (5,88\%), Servicios (68,27\%), Impuestos netos (9,23\%).
    \item \textbf{2023:} Agricultura (2,50\%), Industria (14,72\%), Construcción (5,39\%), Servicios (68,66\%), Impuestos netos (8,72\%).
\end{itemize}

\subsubsection{3. Tasa de crecimiento interanual del PIB (2020-2023)}
Fórmula aplicada: $TV = \left(\frac{PIB_t}{PIB_{t-1}} - 1\right) \times 100$
\begin{itemize}
    \item \textbf{2021:} 9,41\% (respecto a 2020).
    \item \textbf{2022:} 11,18\% (respecto a 2021).
    \item \textbf{2023:} 9,08\% (respecto a 2022).
\end{itemize}

\subsubsection{4. Tasa de variación media del PIB (2008-2023)}
La TVMA para los 15 años del periodo:
\[ \left( \sqrt[15]{\frac{1.498.324}{1.112.432}} - 1 \right) \times 100 = 2,01\% \text{ anual.} \]

\subsubsection{5. TVMA por subperiodos}
\begin{itemize}
    \item \textbf{2008-2014:} -1,13\% (Recesión).
    \item \textbf{2014-2019:} 3,83\% (Expansión).
    \item \textbf{2020-2023:} 9,88\% (Rebote nominal e inflación).
\end{itemize}

\subsubsection{6. PIB en números índices (Base 2020)}
\begin{itemize}
    \item 2008: 98,51 | 2014: 92,01 | 2019: 111,02 | 2020: 100,00 | 2021: 109,41 | 2022: 121,64 | 2023: 132,69.
\end{itemize}

\subsubsection{7. Análisis de fases (2008-2023)}
\begin{description}
    \item[Fase 1 (2008-2014):] Crisis y recesión. Destrucción de valor (-1,13\% anual) y desplome de la construcción.
    \item[Fase 2 (2014-2019):] Recuperación y expansión sólida (3,83\% anual) con fuerte terciarización.
    \item[Fase 3 (2020):] Shock pandémico. Caída brusca del PIB real y nominal.
    \item[Fase 4 (2020-2023):] Rebote nominal intenso marcado por la inflación y la apertura exterior (exportaciones $>$ 38\%).
\end{description}

\subsection{Ejercicio 2: PIB Real y Deflactor}

\textbf{Enunciado:} Con los datos del Cuadro 2, calcular el PIB a precios constantes (2008-2023), el deflactor del PIB, el crecimiento de precios, los componentes a precios constantes y la TVMA nominal vs. real.

\subsubsection{1. PIB a precios constantes (Real)}
Calculado como $(PIB_{base2020} \times \text{Índice Volumen}) / 100$.
\begin{itemize}
    \item 2008: 1.184.545,49 M€ | 2014: 1.102.112,86 M€ | 2019: 1.268.107,32 M€ | 2020: 1.129.214 M€ | 2023: 1.313.275,88 M€.
\end{itemize}

\subsubsection{2. Deflactor del PIB}
Calculado como $(PIB_{corriente} / PIB_{constante}) \times 100$.
\begin{itemize}
    \item 2008: 93,91 | 2014: 94,27 | 2019: 98,86 | 2020: 100,00 | 2023: 114,09.
\end{itemize}

\subsubsection{3. Crecimiento de precios (Inflación)}
\begin{itemize}
    \item \textbf{2008-2014:} 0,38\% (Estancamiento).
    \item \textbf{2014-2019:} 4,87\% (Crecimiento moderado).
    \item \textbf{2020-2023:} 14,09\% (Alta inflación).
\end{itemize}

\subsubsection{4. Componentes a precios constantes (2008 vs 2023)}
\textbf{Demanda (M€):}
\begin{itemize}
    \item Consumo final: 933.685,97 $\rightarrow$ 993.470,01.
    \item FBC: 301.397,19 $\rightarrow$ 265.294,74.
    \item Exportaciones: 308.399,62 $\rightarrow$ 458.469,07.
\end{itemize}

\subsubsection{5. TVMA Nominal vs Real (2008-2023)}
El PIB total creció un 2,01\% nominal frente a un 0,69\% real. El componente más inflacionista ha sido la \textbf{Agricultura} (+2,29 puntos de diferencia) y los \textbf{Impuestos netos}.

\subsection{Ejercicio 3: Crecimiento Trimestral y Estructura}

\subsubsection{1. Crecimiento interanual (T3 2023)}
\begin{itemize}
    \item Consumo final: 2,62\% | FBC: -1,82\% | Exportaciones: 0,08\% | Importaciones: -1,28\%.
    \item Servicios: 2,99\% | Agricultura: 12,45\% (Efecto base tras caída en 2022).
\end{itemize}

\subsubsection{2. Evolución trimestral del PIB (2022-2023)}
Tasas intertrimestrales: T2'22: 1,71\% | T3'22: 0,88\% | T4'22: 0,53\% | T1'23: 0,70\% | T2'23: 0,26\% | T3'23: 0,60\%.

\subsubsection{3. Estructuras porcentuales (Constantes)}
\begin{itemize}
    \item \textbf{Demanda T4 2022:} Consumo (74,80\%), FBC (20,54\%), Exportaciones (35,26\%).
    \item \textbf{Oferta T3 2023:} Servicios (70,86\%), Industria (13,48\%), Construcción (5,30\%).
\end{itemize}

\subsubsection{4. Análisis de contribución}
Los \textbf{Servicios} y el \textbf{Gasto en consumo final} son los motores del crecimiento en 2023 por su elevado peso estructural y dinamismo positivo, compensando el estancamiento industrial y del sector exterior.

\subsection{Ejercicio 4: Crecimiento y Deflactor (2005-2019)}

\subsubsection{1. Comparativa de tasas (Corrientes vs. Constantes)}
La diferencia entre el crecimiento nominal (corriente) y el real (constante) se debe exclusivamente al \textbf{efecto de los precios (inflación)}. Cuando la tasa nominal supera la real (ej. 2006, 2017), existe inflación; si es inferior, deflación.

\subsubsection{2. Deflactor implícito del PIB (Base 2015)}
Representa el nivel general de precios de todos los bienes y servicios producidos.
\begin{itemize}
    \item 2005: 90,15 | 2010: 99,42 | 2015: 100,00 | 2019: 104,23.
\end{itemize}

\subsubsection{3. Fases económicas}
\begin{itemize}
    \item \textbf{Expansión (2005-2008):} Crecimiento real alto (~4\%) con inflación elevada.
    \item \textbf{Crisis (2009-2014):} Recesión y devaluación interna (inflación negativa en 2011, 2012, 2014).
    \item \textbf{Recuperación (2015-2019):} Crecimiento real sólido (~3\%) con inflación moderada y estable (~1,3\%).
\end{itemize}

\subsection{Ejercicio 6: PIB en Valores Monetarios}

\subsubsection{1. Cálculo del Deflactor}
Fórmula: $(\text{Índice PIB Corrientes} / \text{Índice PIB Volumen}) \times 100$.
\begin{itemize}
    \item 1995: 64,31 | 2005: 90,15 | 2020: 105,38.
\end{itemize}

\subsubsection{2. Crecimiento de precios}
\begin{itemize}
    \item Entre 1995 y 2005: 40,18\%.
    \item Entre 2005 y 2020: 16,89\%.
\end{itemize}

\subsubsection{3. Conversión a valores monetarios (PIB Constante)}
Es posible calcularlo usando el valor monetario del año base 2015 (1.077.590 M€).
\begin{itemize}
    \item 1995: 716.166,31 M€ | 2005: 1.028.667,41 M€ | 2020: 1.064.658,92 M€.
\end{itemize}

\subsection{Ejercicio 7: Análisis del IPC (2019-2021)}

\subsubsection{1. Tasas de Inflación}
\begin{itemize}
    \item \textbf{Interanual:} Agosto 2020: -0,52\% | Agosto 2021: 3,30\%.
    \item \textbf{Acumulada (Agosto 2021):} 2,23\%.
    \item \textbf{Media anual 2020:} -0,32\%.
\end{itemize}

\subsubsection{2. Dinámica de precios}
\begin{itemize}
    \item En \textbf{Enero} coinciden las tasas mensual y acumulada por compartir la base de diciembre.
    \item En \textbf{Diciembre} coinciden las tasas interanual y acumulada.
\end{itemize}

\subsection{Ejercicio 8: Deflactación con IPC}

Al deflactar el PIB nominal con el IPC (Base 2021), se observa que en 2022 gran parte del crecimiento nominal fue un \textbf{espejismo inflacionario} (IPC subió 8,40\%), dejando el crecimiento real en un moderado 2,57\%.

\subsection{Ejercicio 9 y 10: Contribución al Crecimiento (2019)}

\begin{itemize}
    \item \textbf{Demanda:} El principal motor fue el \textbf{consumo final} (+1,98 pp), seguido de la FBC (+1,13 pp). El sector exterior aportó +0,32 pp netos.
    \item \textbf{Oferta:} El crecimiento del 3,42\% se sustentó principalmente en los \textbf{Servicios}, que aportaron la mayor parte de la tracción económica.
\end{itemize}

\section{Cuaderno 2: Productividad, Convergencia y Estructura Sectorial}

\subsection{Ejercicio 1: Crecimiento Intensivo y Eficiente}

\textbf{Enunciado:} Estimar la TVMAA de la producción, empleo, población, renta por habitante, productividad aparente y tasa global de empleo entre el Periodo 1 y 10.

\subsubsection{1. Tasas de Crecimiento (TVMAA, n=9)}
\begin{itemize}
    \item \textbf{Producción (PIB):} 4,61\%.
    \item \textbf{Empleo:} 2,05\% | \textbf{Población:} 1,58\%.
    \item \textbf{Renta por habitante:} 2,98\%.
    \item \textbf{Productividad aparente:} 2,51\%.
    \item \textbf{Tasa Global de Empleo (TGE):} 0,47\%.
\end{itemize}

\subsubsection{2. Interpretación}
El crecimiento es sano y eficiente: la producción crece muy por encima del empleo, lo que se traduce en una mejora de la productividad (2,51\%). El nivel de vida (renta por habitante) mejora gracias a la eficiencia tecnológica y, en menor medida, a la mayor ocupación.

\subsection{Ejercicio 2: Contribución a la Renta por Habitante}

Usando la identidad: \textit{TVMAA Renta $\approx$ TVMAA Productividad + TVMAA TGE}.
\begin{itemize}
    \item \textbf{Productividad:} 84,23\% de la mejora de la renta.
    \item \textbf{Tasa de Empleo:} 15,77\% de la mejora.
\end{itemize}
Esto refleja un \textbf{crecimiento intensivo}, ideal para economías desarrolladas donde el factor trabajo tiene límites físicos.

\subsection{Ejercicio 3: Economía Española (2005-2014)}

\subsubsection{1. Evolución de componentes}
\begin{itemize}
    \item \textbf{Renta por habitante:} -0,32\% (Pérdida de bienestar).
    \item \textbf{Productividad Aparente:} +1,29\% (Espejismo por destrucción de empleo de baja productividad).
    \item \textbf{Tasa de Ocupación:} -2,03\% (Causa real del empobrecimiento).
    \item \textbf{Tasa de Actividad:} +0,45\% (Efecto del trabajador añadido).
\end{itemize}

\subsection{Ejercicio 4: Convergencia España-Alemania}

Basado en las tasas históricas (2014-2023): España (3,67\%) y Alemania (3,33\%).
\textbf{Resultado:} España tardaría aproximadamente \textbf{143 años} en alcanzar a Alemania. La convergencia es extremadamente lenta debido al escaso diferencial de crecimiento (0,34 pp) frente a la gran brecha inicial.

\subsection{Ejercicio 5: Ritmo para Convergencia en 15 años}

Para igualar la renta de Alemania en 15 años, España necesitaría una \textbf{TVMAA del 6,61\%}. Este ritmo es utópico para una economía madura, propio de "milagros económicos" en países emergentes.

\subsection{Ejercicio 6: Estructura y Contribución Sectorial}

\subsubsection{1. Estructura del PIB}
\begin{itemize}
    \item Servicios (57,4\%), Industria (17,1\%), Construcción (7,3\%), Agricultura (3,4\%).
\end{itemize}

\subsubsection{2. Contribución al crecimiento}
Los \textbf{Servicios} aportaron 1,90 pp de los 2,80 pp de crecimiento total, consolidándose como el motor absoluto de la economía. La industria sufrió una ligera desindustrialización relativa.

\subsection{Ejercicio 7: VAB y Consumos Intermedios}

\subsubsection{1. Estructura y Dinamismo}
Entre 2006 y 2023, la \textbf{Construcción} colapsó (pasó del 12\% al 6\%), mientras que los \textbf{Servicios} se expandieron hasta superar el 75\% del VAB. Las ramas más dinámicas fueron Energía eléctrica (4,96\%) y Actividades Inmobiliarias (4,09\%).

\subsubsection{2. Consumos Intermedios (CI)}
La industria manufacturera es la que más CI requiere (76\% de su producción). A nivel nacional, los CI bajaron del 55\% al 51\% debido al \textbf{efecto composición}: el mayor peso de los servicios (intensivos en capital humano, no en materias primas) reduce la media de gasto en CI.

\subsection{Ejercicio 8: Distribución de la Renta}

\begin{itemize}
    \item \textbf{Agricultura y Hostelería:} Predomina el EBE/Rentas Mixtas ($>$60-80\%), indicando una estructura de autónomos y pequeños negocios familiares.
    \item \textbf{Industria y Comercio:} Predomina la Remuneración de Asalariados, reflejando empresas de mayor tamaño y empleo formal.
    \item \textbf{Subvenciones:} La Agricultura destaca por sus impuestos netos negativos (-5.339 M€), evidenciando el peso de las ayudas de la PAC.
\end{itemize}

\subsection{Ejercicio 9: Productividad en la Industria (PAT)}

\begin{itemize}
    \item \textbf{Mayor PAT (2022):} Energía (1,15 M€) e Industrias extractivas (392.202 €) por su alta intensidad de capital.
    \item \textbf{Mejor Evolución:} Fabricación de vehículos (+5,55\% anual), reflejando inversión en robótica y automatización.
    \item \textbf{Sectores estancados:} Alimentación y Papel, con crecimientos de eficiencia nulos o negativos.
\end{itemize}

\subsection{Ejercicio 10: Productividad en los Servicios}

La productividad en servicios es, en general, inferior a la industrial. Sin embargo, ramas de \textbf{Servicios Avanzados} (Finanzas, Inmobiliarias) presentan niveles de PAT muy superiores a los servicios tradicionales (Hostelería, Comercio).

\subsection{Ejercicio 11: Relaciones Comerciales}

\begin{itemize}
    \item \textbf{Agricultura:} Superávit comercial y vocación exportadora (produce 100, consume 80).
    \item \textbf{Industria:} Alta dependencia exterior (penetración de importaciones del 50\%).
    \item \textbf{Servicios:} Sector mayoritariamente cerrado con baja penetración (10\%).
\end{itemize}

\subsection{Ejercicio 12 y 13: Precios y Costes Laborales (CLU)}

Se analiza la evolución del deflactor y la contribución de rentas del trabajo y excedentes al aumento de precios. El \textbf{CLU real} permite evaluar la competitividad-precio de la economía: un aumento del CLU indica que los salarios crecen por encima de la productividad, tensionando los precios interiores.

\section{Cuaderno 3: Sector Exterior y Balanza de Pagos}

\subsection{Ejercicio 1: Indicadores de Apertura}

Fórmulas fundamentales para el análisis exterior:
\begin{itemize}
    \item \textbf{Propensión a Exportar/Importar:} $(X \text{ o } M / PIB) \times 100$.
    \item \textbf{Grado de Apertura:} $((X + M) / PIB) \times 100$.
    \item \textbf{Tasa de Cobertura:} $(X / M) \times 100$.
\end{itemize}

\subsection{Ejercicio 2: Evolución de las Economías Europeas (2007-2017)}

\begin{itemize}
    \item \textbf{Internacionalización:} Aumento generalizado del grado de apertura en todos los países.
    \item \textbf{Corrección de desequilibrios:} Países como España y Portugal pasaron de déficits graves en 2007 a superávits comerciales en 2017 gracias a la reorientación exportadora tras la crisis.
    \item \textbf{Potencia exportadora:} Alemania mantiene una posición estructural de fuerte superávit con una tasa de cobertura constante $>$ 115\%.
\end{itemize}

\subsection{Ejercicio 3: Balanza de Pagos y Financiación}

\subsubsection{1. Cuenta Corriente y de Capital}
El saldo conjunto determina la \textbf{Capacidad o Necesidad de Financiación} de la economía. En el caso analizado, el déficit comercial (-88.459 M€) es el principal responsable de la necesidad de financiación externa (-9,75\% del PIB).

\subsubsection{2. Servicios}
La balanza de servicios suele ser positiva en España (Tasa de cobertura del 129\%), compensando parcialmente el déficit de mercancías.

\subsection{Ejercicio 4: España 2017}
\begin{itemize}
    \item \textbf{Cuenta Corriente:} Saldo positivo (+21.511 M€) impulsado por el turismo (cobertura 305\%).
    \item \textbf{Capacidad de Financiación:} +24.194 M€. España fue financiador neto del resto del mundo.
\end{itemize}

\subsection{Ejercicio 5: Cuenta Financiera (MBP6)}
\begin{itemize}
    \item \textbf{Saldo = VNA - VNP.} Un saldo positivo indica que España financia al mundo (ej. Inversiones Directas, Reservas). Un saldo negativo indica financiación del exterior (ej. Otras inversiones).
    \item \textbf{Resultado:} Capacidad de financiación neta de 18.876 M€ tras ajustar errores y omisiones.
\end{itemize}

\subsection{Ejercicio 6: Posición de Inversión Internacional (PII)}
España mantiene una \textbf{posición deudora neta} (PII negativa: -769.662 M€ en 2017), aunque el saldo ha mejorado desde 2007 (-901.736 M€) debido al crecimiento de activos españoles en el exterior.

\section{Cuaderno 4: Mercado de Trabajo y Sector Público}

\subsection{Ejercicio 1: Tasas del Mercado Laboral (2011)}

\subsubsection{1. Cálculo de Tasas}
\begin{itemize}
    \item \textbf{Tasa de Actividad:} 60,01\% (Varones: 67,45\% | Mujeres: 52,92\%).
    \item \textbf{Tasa de Paro:} 21,64\% (Varones: 21,21\% | Mujeres: 22,16\%).
    \item \textbf{Tasa de Empleo:} 47,03\% (Varones: 53,14\% | Mujeres: 41,19\%).
\end{itemize}

\subsubsection{2. Dinámica del Desempleo}
El paro puede subir por \textbf{destrucción de empleo} o por aumento de la \textbf{población activa} (gente que sale a buscar trabajo). Inversamente, baja por \textbf{creación de empleo} o por reducción de la población activa (emigración o efecto desánimo/inactividad).

\subsection{Ejercicio 2: Ciclos Históricos (1986-2011)}

\begin{description}
    \item[1986-1995:] Estancamiento y crisis. Paro roza el 23\%.
    \item[1995-2007:] Boom del empleo. El paro cae al 8,26\% por la incorporación masiva de mujeres e inmigrantes.
    \item[2007-2011:] Gran Recesión. Colapso del mercado laboral, el paro vuelve al 21,64\%.
\end{description}
El mercado español es altamente \textbf{procíclico} y volátil, dependiente del ciclo económico.

\subsection{Ejercicio 3: EPA y Escenario Hipotético}

\subsubsection{1. Evolución 2014-2015}
Fuerte recuperación con la creación de 521.000 empleos, bajando el paro del 24,4\% al 22,0\%.

\subsubsection{2. Escenario 20\% de Paro}
Para que el paro baje al 20\% manteniendo la ocupación de 2015, la tasa de actividad debería caer al \textbf{58,01\%}. Esto demuestra que, sin creación de empleo, la única vía estadística para bajar la tasa de paro es la salida de activos del mercado laboral.

\subsection{Ejercicio 4: Población Desempleada y Tasa de Actividad}
Si en 2011 se hubiera mantenido la tasa de actividad de una década antes (2001), la tasa de paro habría sido del \textbf{11,24\%} en lugar del 21,64\%. Esto subraya el impacto del crecimiento de la población activa en las cifras de desempleo.

\subsection{Ejercicio 6: Presión Fiscal (2013-2017)}

\subsubsection{1. Evolución}
La presión fiscal pasó del 33,83\% al \textbf{34,30\%}, indicando que el Estado recauda una porción mayor de la riqueza nacional.
\subsubsection{2. Composición}
Las \textbf{cotizaciones sociales} son la partida más voluminosa (~36\%), seguidas del IVA y los impuestos directos (IRPF/Sociedades). El crecimiento de la recaudación se apoyó fuertemente en la recuperación del consumo (IVA).
